% Generated draft: requires mathematical and source review.
\documentclass[11pt]{article}
\usepackage[margin=25mm]{geometry}
\usepackage{fontspec}
\setmainfont{Latin Modern Roman}
\usepackage{amsmath,amssymb,mathtools,unicode-math}
\setmathfont{Latin Modern Math}
\usepackage{enumitem,needspace,xurl,hyperref}
\setlength{\parindent}{0pt}
\setlength{\parskip}{5pt}
\setlength{\emergencystretch}{3em}
\newcommand{\N}{\mathbb{N}}
\newcommand{\Z}{\mathbb{Z}}
\newcommand{\Q}{\mathbb{Q}}
\newcommand{\R}{\mathbb{R}}
\newcommand{\C}{\mathbb{C}}
\newcommand{\F}{\mathbb{F}}
\newcommand{\A}{\mathbb{A}}
\newcommand{\PP}{\mathbb{P}}
\newcommand{\E}{\mathbb{E}}
\newcommand{\eps}{\varepsilon}
\newcommand{\dd}{\,\mathrm{d}}
\begin{document}
\section*{UnsolvedMath ID: 700696}
\textbf{Source problem code:} catalogue entry
\subsection{Definitions and mathematical statement}
Under the standard finite-group interpretation, let \(G\) be a finite group and let \(p\) be a prime. For a subgroup \(X\leq G\), define its normalizer by
\[
N_G(X)=\{g\in G:gXg^{-1}=X\}.
\]
A \(p\)-subgroup is a subgroup whose order is a power of \(p\). A subgroup \(H\leq G\) is \(p\)-local if \(H=N_G(P)\) for some nonidentity \(p\)-subgroup \(P\leq G\). Write \(O_p(H)\) for the largest normal \(p\)-subgroup of \(H\). Following the supplied statement, a \(p\)-local subgroup \(H\) is called \emph{superlocal} if
\[
H=N_G\bigl(O_p(H)\bigr).
\]
\textbf{Problem.} Describe, for each prime \(p\), all superlocal subgroups in alternating groups and in groups of Lie type; that is, determine
\[
\bigl\{H\leq G:\ H=N_G(P)\text{ for some nonidentity }p\text{-subgroup }P\leq G,\quad H=N_G(O_p(H))\bigr\}
\]
for groups \(G\) in the stated families.

\textbf{Unresolved scope qualifications.} The definitions above use the conventional finite-group meaning of \(p\)-local. The supplied source does not explicitly state finiteness, specify the degrees of the alternating groups, or identify the intended forms and parameters of the groups of Lie type. In particular, it does not specify whether the latter family comprises finite simple groups, particular covering or central quotient forms, or some broader class. These scope choices require clarification for a complete classification target. The source also does not specify whether a description up to conjugacy is intended.
\subsection{Short English statement}
For each prime \(p\), describe the \(p\)-local subgroups of alternating groups and groups of Lie type that equal the normalizer of their largest normal \(p\)-subgroup. The finite-group interpretation and the precise scope of the families require human review.
\subsection{Sources}
\begin{itemize}
\item Local UnsolvedMath registry, permanent ID \(700696\), exact source problem code \texttt{catalogue entry}, title \emph{Kourovka Notebook Problem 11.3}; supplied upstream revision \texttt{d371ee917ac9647fbb06cd6fc529deec6f3dda82}.
\item Supplied source attribution: \emph{The Kourovka Notebook (Unsolved Problems in Group Theory)}, 11th issue (1990), Problem 11.3; also identified with the 21st issue edition (2026), arXiv:1401.0300v48. The supplied collection description names E. I. Khukhro and V. D. Mazurov as editors.
\item The registry and background name R. Zh. Aleev as proposer, with proposed year 1990. The statement itself begins with the attribution ``(M. Aschbacher)''. Both attributions are retained; their respective roles are not explained in the supplied record.
\item The supplied registry labels the problem \emph{open}. This is a supplied status label, not an independently verified mathematical or literature assessment. The cited source has not been independently consulted for this draft.
\end{itemize}
\end{document}
